%% ===================================================================
%%  A lower bound on the MSE of RIS-based analog linear synthesis
%%  -- MACRO-FREE VERSION --
%%
%%  Every symbol is written out in full. There are no \newcommand
%%  definitions, so the body can be pasted into any existing project
%%  without copying a preamble.
%%
%%  Requires only: amsmath, amssymb, amsthm  (\boldsymbol comes with
%%  amsmath, so the bm package is not needed).
%%
%%  If your project already defines theorem/lemma/corollary/remark
%%  environments, delete the \newtheorem block below.
%% ===================================================================

\documentclass[11pt]{article}

\usepackage{amsmath,amsfonts,amssymb,amsthm}
\usepackage{geometry}
\geometry{margin=1in}

\newtheorem{theorem}{Theorem}
\newtheorem{lemma}{Lemma}
\newtheorem{corollary}{Corollary}
\theoremstyle{remark}
\newtheorem{remark}{Remark}

\title{A Lower Bound on the MSE of RIS-Based Analog Linear Synthesis}
\author{}
\date{}

\begin{document}
\maketitle

%% =================== BODY STARTS HERE ===================

\section{System model and notation}

A reconfigurable intelligent surface (RIS) of $M$ phase-shifting elements is
asked to synthesize a target latent vector $\mathbf{y}\in\mathbb{C}^{N_r}$ from
an input latent vector $\mathbf{s}\in\mathbb{C}^{N_t}$ during propagation. The
forward model is
\begin{equation}
\mathbf{y}_{\mathrm{synth}}
= \mathbf{H}_2\,\boldsymbol{\Phi}\,\mathbf{H}_1\,\mathbf{s} + \mathbf{w},
\qquad
\boldsymbol{\Phi} = \operatorname{diag}(\boldsymbol{\phi}),
\qquad
\mathbf{w}\sim\mathcal{CN}(\mathbf{0},\sigma^2\mathbf{I}_{N_r}),
\label{eq:forward}
\end{equation}
with $\mathbf{H}_1\in\mathbb{C}^{M\times N_t}$,
$\mathbf{H}_2\in\mathbb{C}^{N_r\times M}$, and the control variable confined to
the complex torus
\begin{equation}
\boldsymbol{\phi}\in\mathbb{T}^{M} \triangleq
\bigl\{\boldsymbol{\phi}\in\mathbb{C}^{M}:\ |\phi_m|=1,\ m=1,\dots,M\bigr\}.
\label{eq:torus}
\end{equation}

Both links are Rician with factor $\kappa$. Writing
$\varepsilon\triangleq 1/(\kappa+1)\in(0,1]$,
\begin{equation}
\mathbf{H}_1=\sqrt{1-\varepsilon}\,\mathbf{a}_{\mathrm{ris},1}
             \mathbf{a}_{\mathrm{tx}}^{H}
            +\sqrt{\varepsilon}\,\tilde{\mathbf{H}}_1,
\qquad
\mathbf{H}_2=\sqrt{1-\varepsilon}\,\mathbf{a}_{\mathrm{rx}}
             \mathbf{a}_{\mathrm{ris},2}^{H}
            +\sqrt{\varepsilon}\,\tilde{\mathbf{H}}_2.
\label{eq:rician}
\end{equation}

\paragraph{Standing assumptions.}
\begin{itemize}
\item[(A1)] $\|\mathbf{a}_{\mathrm{tx}}\|_2^2=N_t$,
      $\|\mathbf{a}_{\mathrm{rx}}\|_2^2=N_r$, and
      $\bigl|[\mathbf{a}_{\mathrm{ris},i}]_m\bigr|=1$ for all $m$ and
      $i\in\{1,2\}$, hence $\|\mathbf{a}_{\mathrm{ris},i}\|_2^2=M$.
\item[(A2)] $\tilde{\mathbf{H}}_1,\tilde{\mathbf{H}}_2$ have i.i.d.\
      $\mathcal{CN}(0,1)$ entries, mutually independent and independent of
      $\mathbf{w}$.
\item[(A3)] $\alpha\triangleq|\mathbf{a}_{\mathrm{tx}}^{H}\mathbf{s}|>0$ and
      $\beta\triangleq\|\mathbf{s}\|_2$.
\item[(A4)] $M\ge N_r\ge 2$.
\end{itemize}

All statements are conditional on a channel realization; probabilistic
statements refer to the draw of $(\tilde{\mathbf{H}}_1,\tilde{\mathbf{H}}_2)$.
Throughout,
\begin{equation}
\mathbf{P}^{\perp} \triangleq
\mathbf{I}_{N_r}-\frac{\mathbf{a}_{\mathrm{rx}}\mathbf{a}_{\mathrm{rx}}^{H}}{N_r}
\label{eq:Pperp}
\end{equation}
is the orthogonal projector onto the complement of the receive beam direction,
and $\bigl(x\bigr)_{+}\triangleq\max(x,0)$. The quantity of interest is
\begin{equation}
\mathrm{MSE}_{\min}(\mathbf{s})
\triangleq
\min_{\boldsymbol{\phi}\in\mathbb{T}^{M}}\
\mathbb{E}_{\mathbf{w}}
\Bigl[\bigl\|\mathbf{y}-(\mathbf{H}_2\boldsymbol{\Phi}\mathbf{H}_1\mathbf{s}
       +\mathbf{w})\bigr\|_2^2\Bigr].
\label{eq:objective}
\end{equation}

\section{Reduction to a linear synthesis problem}

\begin{lemma}[Reduction]
\label{lem:reduction}
With
$\mathbf{A}\equiv\mathbf{A}(\mathbf{s})
 \triangleq\mathbf{H}_2\operatorname{diag}(\mathbf{H}_1\mathbf{s})
 \in\mathbb{C}^{N_r\times M}$,
\begin{equation}
\mathrm{MSE}_{\min}(\mathbf{s})=N_r\sigma^2+\mathcal{E}(\mathbf{s}),
\qquad
\mathcal{E}(\mathbf{s})\triangleq
\min_{\boldsymbol{\phi}\in\mathbb{T}^{M}}\|\mathbf{y}-\mathbf{A}\boldsymbol{\phi}\|_2^2,
\label{eq:reduction}
\end{equation}
and the $m$-th column of $\mathbf{A}$ is
$\mathbf{A}\mathbf{e}_m=[\mathbf{H}_1\mathbf{s}]_m\,\mathbf{h}_{2,m}$, where
$\mathbf{h}_{2,m}$ is the $m$-th column of $\mathbf{H}_2$.
\end{lemma}

\begin{proof}
For deterministic $\mathbf{r}=\mathbf{y}-\mathbf{A}\boldsymbol{\phi}$,
$\|\mathbf{r}-\mathbf{w}\|_2^2
 =\|\mathbf{r}\|_2^2-2\operatorname{Re}(\mathbf{r}^{H}\mathbf{w})
  +\|\mathbf{w}\|_2^2$;
the cross term vanishes since $\mathbb{E}[\mathbf{w}]=\mathbf{0}$, and
$\mathbb{E}\|\mathbf{w}\|_2^2=\operatorname{tr}(\sigma^2\mathbf{I}_{N_r})
 =N_r\sigma^2$.
Entrywise
$[\boldsymbol{\Phi}\mathbf{H}_1\mathbf{s}]_m=[\mathbf{H}_1\mathbf{s}]_m\phi_m$,
so $\boldsymbol{\Phi}\mathbf{H}_1\mathbf{s}
 =\operatorname{diag}(\mathbf{H}_1\mathbf{s})\boldsymbol{\phi}$ and
$\mathbf{H}_2\boldsymbol{\Phi}\mathbf{H}_1\mathbf{s}
 =\mathbf{A}\boldsymbol{\phi}$. The column formula follows from
$\mathbf{A}\mathbf{e}_m=\mathbf{H}_2\mathbf{e}_m[\mathbf{H}_1\mathbf{s}]_m$. The
minimum is attained because $\mathbb{T}^{M}$ is compact and the objective
continuous.
\end{proof}

The additive term $N_r\sigma^2$ is independent of $\boldsymbol{\phi}$ and must be
carried throughout; everything below bounds the \emph{synthesis residual}
$\mathcal{E}(\mathbf{s})$.

\section{The geometric floor}

\begin{theorem}[Range floor; exactness in the pure-LoS limit]
\label{thm:floor}
Let $\mathbf{P}^{\perp}_{\mathbf{A}}$ project onto
$\operatorname{range}(\mathbf{A}(\mathbf{s}))^{\perp}$. Then
\begin{equation}
\mathcal{E}(\mathbf{s})\ \ge\
\inf_{\boldsymbol{\phi}\in\mathbb{C}^{M}}
 \|\mathbf{y}-\mathbf{A}\boldsymbol{\phi}\|_2^2
=\bigl\|\mathbf{P}^{\perp}_{\mathbf{A}}\mathbf{y}\bigr\|_2^2,
\qquad
\operatorname{range}(\mathbf{A})
=\operatorname{span}\bigl\{\mathbf{h}_{2,m}:[\mathbf{H}_1\mathbf{s}]_m\neq0\bigr\}
\subseteq\operatorname{range}(\mathbf{H}_2).
\label{eq:rangefloor}
\end{equation}
In the pure-LoS case $\varepsilon=0$,
\begin{equation}
\mathbf{A}(\mathbf{s})
=(\mathbf{a}_{\mathrm{tx}}^{H}\mathbf{s})\,\mathbf{a}_{\mathrm{rx}}\,
 \mathbf{a}_{\mathrm{ris},2}^{H}
 \operatorname{diag}(\mathbf{a}_{\mathrm{ris},1}),
\qquad
\operatorname{rank}\mathbf{A}=1,
\qquad
\operatorname{range}(\mathbf{A})=\operatorname{span}\{\mathbf{a}_{\mathrm{rx}}\},
\label{eq:rank1}
\end{equation}
so that
\begin{equation}
\mathcal{E}(\mathbf{s})\ \ge\
\bigl\|\mathbf{P}^{\perp}\mathbf{y}\bigr\|_2^2
=\|\mathbf{y}\|_2^2-\frac{|\mathbf{a}_{\mathrm{rx}}^{H}\mathbf{y}|^2}{N_r},
\label{eq:losfloor}
\end{equation}
with \emph{equality} whenever
$|\mathbf{a}_{\mathrm{rx}}^{H}\mathbf{y}|/N_r\le M\alpha$.
\end{theorem}

\begin{proof}
Since $\mathbf{A}\boldsymbol{\phi}\in\operatorname{range}(\mathbf{A})$ for every
$\boldsymbol{\phi}$ and $\mathbb{T}^{M}\subset\mathbb{C}^{M}$, the residual is at
least the distance from $\mathbf{y}$ to $\operatorname{range}(\mathbf{A})$,
attained by orthogonal projection. The range identity follows from
Lemma~\ref{lem:reduction}: the $m$-th column spans the same line as
$\mathbf{h}_{2,m}$ when $[\mathbf{H}_1\mathbf{s}]_m\neq0$ and vanishes otherwise.

For $\varepsilon=0$ we have
$\mathbf{H}_1=\mathbf{a}_{\mathrm{ris},1}\mathbf{a}_{\mathrm{tx}}^{H}$ and
$\mathbf{H}_2=\mathbf{a}_{\mathrm{rx}}\mathbf{a}_{\mathrm{ris},2}^{H}$, hence
$\mathbf{H}_1\mathbf{s}
 =(\mathbf{a}_{\mathrm{tx}}^{H}\mathbf{s})\mathbf{a}_{\mathrm{ris},1}$, whose
entries all have modulus $\alpha>0$ by (A1) and (A3); this gives
\eqref{eq:rank1} and therefore \eqref{eq:losfloor}.

For exactness, \eqref{eq:rank1} yields
\begin{equation}
\mathbf{A}\boldsymbol{\phi}
=(\mathbf{a}_{\mathrm{tx}}^{H}\mathbf{s})\,c(\boldsymbol{\phi})\,
 \mathbf{a}_{\mathrm{rx}},
\qquad
c(\boldsymbol{\phi})=\sum_{m=1}^{M}
 \overline{[\mathbf{a}_{\mathrm{ris},2}]_m}\,
 [\mathbf{a}_{\mathrm{ris},1}]_m\,\phi_m .
\label{eq:cphi}
\end{equation}
Each summand has unit modulus and arbitrary phase as $\phi_m$ ranges over the
unit circle, so $\{c(\boldsymbol{\phi}):\boldsymbol{\phi}\in\mathbb{T}^{M}\}$ is
exactly the closed disk of radius $M$: containment is the triangle inequality,
and conversely, fixing a target phase $\theta$ and pairing elements as
$e^{i(\theta\pm\omega)}$ realizes $2\cos\omega\,e^{i\theta}$ per pair, so every
modulus in $[0,M]$ is attainable (for odd $M$, use $\lfloor M/2\rfloor$ pairs
plus one free element). The reachable set is thus
$\{z\mathbf{a}_{\mathrm{rx}}:|z|\le M\alpha\}$, and the minimizing
$z^{\star}=\mathbf{a}_{\mathrm{rx}}^{H}\mathbf{y}/N_r$ is feasible precisely
under the stated condition.
\end{proof}

\begin{remark}
Under pure LoS the RIS controls exactly one complex scalar: the received vector
is locked to the direction $\mathbf{a}_{\mathrm{rx}}$ irrespective of
$\mathbf{s}$ and $\boldsymbol{\phi}$. Neither the transmit power nor the surface
size $M$ alters this; $M$ only sets the radius of the reachable disk.
\end{remark}

\begin{remark}
At any finite $\kappa$ with $M\ge N_r$, $\mathbf{A}(\mathbf{s})$ has full row
rank almost surely and \eqref{eq:rangefloor} degenerates to $0$. The
finite-$\kappa$ statement therefore cannot be a rank statement; it must be a
conditioning statement, which is Theorem~\ref{thm:finitekappa}.
\end{remark}

\section{Non-asymptotic floor at finite Rician factor}

\begin{theorem}[Finite $\kappa$]
\label{thm:finitekappa}
Deterministically,
\begin{equation}
\mathcal{E}(\mathbf{s})\ \ge\
\Bigl(\bigl\|\mathbf{P}^{\perp}\mathbf{y}\bigr\|_2
 -\sqrt{M}\,\bigl\|\mathbf{P}^{\perp}\mathbf{A}(\mathbf{s})\bigr\|_2
\Bigr)_{+}^{2},
\qquad
\bigl\|\mathbf{P}^{\perp}\mathbf{A}(\mathbf{s})\bigr\|_2
\le\sqrt{\varepsilon}\,s_{\max}(\tilde{\mathbf{H}}_2)\,
   \max_{m}\bigl|[\mathbf{H}_1\mathbf{s}]_m\bigr|.
\label{eq:finitekappa}
\end{equation}
Consequently, for any $\delta\in(0,1)$, with probability at least $1-\delta$
over the NLoS draw,
\begin{equation}
\mathrm{MSE}_{\min}(\mathbf{s})\ \ge\ N_r\sigma^2
+\Bigl(\bigl\|\mathbf{P}^{\perp}\mathbf{y}\bigr\|_2
 -\Delta(\kappa,M,\delta)\Bigr)_{+}^{2},
\label{eq:finitekappaprob}
\end{equation}
where, with $\tau=\sqrt{2\ln(2/\delta)}$,
\begin{equation}
\Delta(\kappa,M,\delta)
=\sqrt{\frac{M}{\kappa+1}}
 \Bigl(\sqrt{M}+\sqrt{N_r}+\tau\Bigr)
 \Bigl(\alpha+\tfrac{\beta}{\sqrt{\kappa+1}}
  \sqrt{\ln\tfrac{2M}{\delta}}\Bigr)
=\mathcal{O}\!\left(\frac{M\alpha}{\sqrt{\kappa}}\right).
\label{eq:delta}
\end{equation}
In particular
$\liminf_{\kappa\to\infty}\mathcal{E}(\mathbf{s})
 \ge\|\mathbf{P}^{\perp}\mathbf{y}\|_2^2$, recovering
Theorem~\ref{thm:floor} continuously. The bound is non-vacuous as soon as
$\kappa+1\gtrsim M^2\alpha^2/\|\mathbf{P}^{\perp}\mathbf{y}\|_2^2$.
\end{theorem}

\begin{proof}
By (A1), $\mathbf{P}^{\perp}\mathbf{a}_{\mathrm{rx}}=\mathbf{0}$, so the
rank-one LoS term of $\mathbf{H}_2$ is annihilated exactly and
\begin{equation}
\mathbf{P}^{\perp}\mathbf{H}_2
=\sqrt{\varepsilon}\,\mathbf{P}^{\perp}\tilde{\mathbf{H}}_2 .
\label{eq:annihilation}
\end{equation}
Fix any $\boldsymbol{\phi}\in\mathbb{T}^{M}$. Since $\mathbf{P}^{\perp}$ is an
orthogonal projector, $\|\mathbf{r}\|_2\ge\|\mathbf{P}^{\perp}\mathbf{r}\|_2$,
whence by the reverse triangle inequality and the definition of the spectral
norm
\begin{equation}
\|\mathbf{y}-\mathbf{A}\boldsymbol{\phi}\|_2
\ \ge\ \bigl\|\mathbf{P}^{\perp}\mathbf{y}
        -\mathbf{P}^{\perp}\mathbf{A}\boldsymbol{\phi}\bigr\|_2
\ \ge\ \bigl\|\mathbf{P}^{\perp}\mathbf{y}\bigr\|_2
       -\bigl\|\mathbf{P}^{\perp}\mathbf{A}\bigr\|_2\|\boldsymbol{\phi}\|_2 .
\label{eq:revtriangle}
\end{equation}
On the torus $\|\boldsymbol{\phi}\|_2=\sqrt{M}$ exactly; taking positive parts,
squaring and minimizing over $\boldsymbol{\phi}$ gives the first part of
\eqref{eq:finitekappa}. For the second part,
$\mathbf{P}^{\perp}\mathbf{A}
 =\mathbf{P}^{\perp}\mathbf{H}_2\operatorname{diag}(\mathbf{H}_1\mathbf{s})$
with
$\|\operatorname{diag}(\mathbf{H}_1\mathbf{s})\|_2
 =\max_m|[\mathbf{H}_1\mathbf{s}]_m|$;
apply submultiplicativity and \eqref{eq:annihilation}.

For \eqref{eq:finitekappaprob}, the Davidson--Szarek bound gives
$\Pr\bigl(s_{\max}(\tilde{\mathbf{H}}_2)>\sqrt{M}+\sqrt{N_r}+\tau\bigr)
 \le e^{-\tau^2/2}$,
which is at most $\delta/2$ for $\tau=\sqrt{2\ln(2/\delta)}$. Writing
$z_m\triangleq\tilde{\mathbf{g}}_{1,m}^{T}\mathbf{s}\sim\mathcal{CN}(0,\beta^2)$,
the variable $|z_m|/\beta$ is Rayleigh with
$\Pr(|z_m|/\beta>t)=e^{-t^2}$, so a union bound over $m$ gives
$\max_m|z_m|\le\beta\sqrt{\ln(2M/\delta)}$ with probability at least
$1-\delta/2$; combined with
$|[\mathbf{H}_1\mathbf{s}]_m|
 \le\sqrt{1-\varepsilon}\,\alpha+\sqrt{\varepsilon}|z_m|$
and a final union bound, \eqref{eq:delta} follows. Add $N_r\sigma^2$ via
\eqref{eq:reduction}. The rate holds since
$\sqrt{M}(\sqrt{M}+\sqrt{N_r}+\tau)=\Theta(M)$ for $M\ge N_r$.
\end{proof}

\section{Consequences}

\begin{corollary}[Cosine objective with automatic gain control]
\label{cor:cosine}
For the scale-invariant matching objective,
\begin{equation}
\max_{\boldsymbol{\phi}\in\mathbb{T}^{M}}
\frac{\bigl|\langle\mathbf{y},\mathbf{A}\boldsymbol{\phi}\rangle\bigr|}
     {\|\mathbf{y}\|_2\,\|\mathbf{A}\boldsymbol{\phi}\|_2}
\ \le\
\frac{\|\mathbf{P}_{\mathbf{A}}\mathbf{y}\|_2}{\|\mathbf{y}\|_2}
\ \xrightarrow[\ \varepsilon\to0\ ]{}\
\frac{|\mathbf{a}_{\mathrm{rx}}^{H}\mathbf{y}|}
     {\sqrt{N_r}\,\|\mathbf{y}\|_2},
\label{eq:cosbound}
\end{equation}
and the gain-matched squared error obeys the same floor,
\begin{equation}
\min_{c>0}\|\mathbf{y}-c\,\mathbf{A}\boldsymbol{\phi}\|_2^2
=\|\mathbf{y}\|_2^2\bigl(1-\cos^2\theta\bigr)
\ \ge\ \bigl\|\mathbf{P}^{\perp}_{\mathbf{A}}\mathbf{y}\bigr\|_2^2 .
\label{eq:agc}
\end{equation}
\end{corollary}

\begin{proof}
Cauchy--Schwarz applied to
$\langle\mathbf{y},\mathbf{v}\rangle
 =\langle\mathbf{P}_{\mathbf{A}}\mathbf{y},\mathbf{v}\rangle$ for
$\mathbf{v}\in\operatorname{range}(\mathbf{A})$, together with the
scale-optimized residual identity
$\min_c\|\mathbf{y}-c\hat{\mathbf{v}}\|_2^2
 =\|\mathbf{y}\|_2^2-|\hat{\mathbf{v}}^{H}\mathbf{y}|^2$ for unit
$\hat{\mathbf{v}}$.
\end{proof}

\begin{corollary}[Digital pre- and post-processing does not escape the floor]
\label{cor:airfc}
Let
$\mathbf{W}_{\mathrm{phys}}
 =\mathbf{U}^{H}\mathbf{H}_2\boldsymbol{\Phi}\mathbf{H}_1
  \mathbf{P}_{\mathrm{tx}}$
for arbitrary $\mathbf{P}_{\mathrm{tx}}\in\mathbb{C}^{N_t\times N_t}$,
$\mathbf{U}\in\mathbb{C}^{N_r\times N_r}$ and
$\boldsymbol{\phi}\in\mathbb{T}^{M}$. At $\varepsilon=0$,
$\operatorname{rank}(\mathbf{W}_{\mathrm{phys}})\le1$, and therefore for any
target $\mathbf{W}$,
\begin{equation}
\frac{\|\mathbf{W}_{\mathrm{phys}}-\mathbf{W}\|_F}{\|\mathbf{W}\|_F}
\ \ge\
\sqrt{1-\frac{\sigma_1^2(\mathbf{W})}{\|\mathbf{W}\|_F^2}} .
\label{eq:eckartyoung}
\end{equation}
\end{corollary}

\begin{proof}
By \eqref{eq:rank1},
$\mathbf{H}_2\boldsymbol{\Phi}\mathbf{H}_1
 =c(\boldsymbol{\phi})\mathbf{a}_{\mathrm{rx}}\mathbf{a}_{\mathrm{tx}}^{H}$
is rank one for every $\boldsymbol{\phi}$, and rank does not increase under left
or right multiplication. Apply the Eckart--Young theorem in the Frobenius norm.
\end{proof}

\section{The spectral route, for reference}

Let
$\mathbf{R}_{\mathrm{eq}}\triangleq\mathbf{A}\mathbf{A}^{H}
 =\sum_{k=1}^{N_r}\lambda_k\mathbf{u}_k\mathbf{u}_k^{H}$
with $\lambda_1\ge\cdots\ge\lambda_{N_r}\ge0$, and
$y_k\triangleq\mathbf{u}_k^{H}\mathbf{y}$. Relaxing $\mathbb{T}^{M}$ to the ball
$\{\|\boldsymbol{\phi}\|_2^2\le M\}$ and invoking weak duality for the resulting
QCQP yields, \emph{unconditionally},
\begin{equation}
\mathcal{E}(\mathbf{s})\ \ge\ \sup_{\mu\ge0}
\left[\sum_{k=1}^{N_r}\frac{\mu\,|y_k|^2}{\lambda_k+\mu}-\mu M\right],
\label{eq:dual}
\end{equation}
with the supremum equal to the exact value of the relaxed problem by Slater's
condition. Under \eqref{eq:rician}, with probability at least $1-\delta$ and for
$M\ge N_r\ge3$,
\begin{equation}
\lambda_1=\Bigl(\tfrac{\kappa}{\kappa+1}\Bigr)^{2}MN_r\alpha^2
          \bigl(1+o(1)\bigr),
\qquad
\lambda_k=\frac{c_k\,\kappa\,M\,\alpha^2}{(\kappa+1)^2},
\quad 2\le k\le N_r-1,
\label{eq:eigs}
\end{equation}
with $c_k=\Theta(1)$ determined by the Wishart edge. Dropping the leading mode
and bounding $\lambda_{k\ge2}\le\bar\lambda$ in \eqref{eq:dual} gives the closed
form
\begin{equation}
\sup_{\mu\ge0}\left[\frac{\mu Y}{\bar\lambda+\mu}-\mu M\right]
=\Bigl(\sqrt{Y}-\sqrt{\bar\lambda M}\Bigr)_{+}^{2}
=\Bigl(\sqrt{Y}-\Theta\bigl(M\alpha/\sqrt{\kappa}\bigr)\Bigr)_{+}^{2},
\qquad
Y\triangleq\sum_{k\ge2}|y_k|^2,
\label{eq:closedform}
\end{equation}
attained at $\mu=\sqrt{Y\bar\lambda/M}-\bar\lambda$. This matches the
$\mathcal{O}(M\alpha/\sqrt{\kappa})$ rate of Theorem~\ref{thm:finitekappa},
confirming that the two routes agree.

\begin{remark}[Validity of the squared-denominator form]
\label{rem:errata}
Let
$G(\mu)=\sum_k\bigl(\tfrac{\mu}{\lambda_k+\mu}\bigr)^2|y_k|^2$ and
$h(\mu)=\sum_k\lambda_k|y_k|^2/(\lambda_k+\mu)^2$. Then $G(\sigma^2)$
lower-bounds $\mathcal{E}(\mathbf{s})$ \emph{if and only if}
$h(\sigma^2)\ge M$, i.e.\ iff the KKT multiplier $\mu^{\star}$ of the power
constraint (pinned by $h(\mu^{\star})=M$) satisfies $\mu^{\star}\ge\sigma^2$.
This condition holds only in the intermediate band
\begin{equation}
\frac{M\alpha^2}{Y}\ \lesssim\ \kappa\ \lesssim\ \frac{M\alpha^2}{\sigma^2},
\label{eq:band}
\end{equation}
and fails at both ends, in particular as $\kappa\to\infty$. Identifying the KKT
multiplier with a Tikhonov regularization parameter $\mu\propto\sigma^2$ is
therefore not a valid proof step; \eqref{eq:dual} is the correct unconditional
substitute, and carries exponent $1$ together with the penalty $-\mu M$.
\end{remark}

\section{Assembled result}

Collecting \eqref{eq:reduction}, \eqref{eq:dual}, \eqref{eq:finitekappaprob} and
\eqref{eq:eigs}, with probability at least $1-\delta$,
\begin{equation}
\mathrm{MSE}_{\min}(\mathbf{s})\ \ge\ N_r\sigma^2+\max\left\{
\underbrace{\Bigl(\|\mathbf{P}^{\perp}\mathbf{y}\|_2
            -\Delta(\kappa,M,\delta)\Bigr)_{+}^{2}}_{\text{geometric}},
\ \
\underbrace{\sup_{\mu\ge0}\left[\sum_{k}
 \frac{\mu\,|\mathbf{u}_k^{H}\mathbf{y}|^2}{\lambda_k+\mu}-\mu M\right]}
 _{\text{spectral}}
\right\}.
\label{eq:assembled}
\end{equation}
Taking $\kappa\to\infty$ at fixed $M,N_r,\sigma^2$ gives $\Delta\to0$ and
$\mathbf{u}_1\to\mathbf{a}_{\mathrm{rx}}/\sqrt{N_r}$ by Davis--Kahan, since the
spectral gap $\lambda_1-\lambda_2=\Theta(MN_r\alpha^2)$ dominates the LoS--NLoS
coupling $\mathcal{O}(\sqrt{\varepsilon M}N_r)$. Hence
\begin{equation}
\lim_{\kappa\to\infty}\mathrm{MSE}_{\min}(\mathbf{s})
\ \ge\ N_r\sigma^2
+\left\|\Bigl(\mathbf{I}_{N_r}
 -\tfrac{\mathbf{a}_{\mathrm{rx}}\mathbf{a}_{\mathrm{rx}}^{H}}{N_r}\Bigr)
 \mathbf{y}\right\|_2^2
= N_r\sigma^2+\|\mathbf{y}\|_2^2-|\mathbf{u}_1^{H}\mathbf{y}|^2 .
\label{eq:limit}
\end{equation}
No secondary limit $\mathrm{SNR}\to\infty$ is required: the geometric floor is a
statement about $\operatorname{range}(\mathbf{A}(\mathbf{s}))$ and is
independent of $\sigma^2$, which enters only through the separate additive term
$N_r\sigma^2$.

\begin{remark}[Consequence]
If the target is $\mathbf{y}=\mathbf{W}_{\mathrm{lin}}\mathbf{s}$ for a
full-rank $\mathbf{W}_{\mathrm{lin}}\in\mathbb{C}^{N_r\times N_t}$, and
$\mathbf{s}$ is not concentrated on
$\mathbf{W}_{\mathrm{lin}}^{-1}\operatorname{span}
 \{\mathbf{a}_{\mathrm{rx}}\}$, then for $N_r\ge2$
\begin{equation}
\mathbb{E}_{\mathbf{s}}
 \bigl\|\mathbf{P}^{\perp}\mathbf{W}_{\mathrm{lin}}\mathbf{s}\bigr\|_2^2
=\Theta\bigl(\mathbb{E}_{\mathbf{s}}
 \|\mathbf{W}_{\mathrm{lin}}\mathbf{s}\|_2^2\bigr),
\label{eq:consequence}
\end{equation}
i.e.\ a LoS-dominated link cannot carry a full-rank linear map at any SNR and
with any number of RIS elements. Rich scattering is what makes
$\|\mathbf{P}^{\perp}_{\mathbf{A}}\mathbf{y}\|_2=0$ attainable.
\end{remark}

\section{Limitations}

The bounds are one-sided: they constrain what the surface cannot do, and do not
certify that a particular optimizer approaches them. The relaxation underlying
\eqref{eq:dual} is shown to be lossless in the rank-one regime, but its gap at
intermediate $\kappa$ is not quantified here. All statements are per-input and
per-realization; averaging over a data distribution requires
$\mathbf{a}_{\mathrm{tx}}^{H}\mathbf{s}\neq0$ almost surely and integrable
$\beta/\alpha$. The expansion \eqref{eq:eigs} requires
$\kappa\gtrsim\ln M\cdot(\beta/\alpha)^2$, so the regime $M\to\infty$ at fixed
$\kappa$ is not covered. Path loss is omitted; it rescales all $\lambda_k$ and
$\Delta$ by a common factor and leaves \eqref{eq:losfloor} unaffected.

%% =================== BODY ENDS HERE ===================

\end{document}
